\chapter{项目背景与目的}

\section{任务目的}

废塑料瓶回收通常要求瓶身与瓶盖分离，便于后续清洗、材质分类和再利用。本项目希望机器人直接面对桌面上数量、位置和方向不固定的瓶子，完成“取瓶—开盖—分别投放”的连续工作。这里的目标不是只做一个展示动作，而是形成可持续扩充数据、重复训练、长期运行和量化评测的系统。

\plain{人看到瓶子时，会自然地判断“有没有盖、哪一端是瓶口、怎么抓才不会滑”。机器人必须从画面和过去的示教里学会这些判断，还要让两只手在同一时刻做不同但互相配合的动作。}

\section{为什么选择 ALOHA}

ALOHA 的核心价值是“双臂工作 + 双臂示教”：操作者通过前侧两只示教臂自然地演示动作，后侧两只工作臂同步完成操作并记录多视角画面和关节动作。原始 ALOHA 工作证明该平台适合采集精细双臂操作示教；ALOHA 2 又进一步提高了人体工学、稳健性和规模化采集能力\cite{zhao2023aloha,aloha2_2024}。

对本项目而言，选择 ALOHA 有三个直接原因：

\begin{enumerate}
  \item 一只工作臂可以稳定瓶身，另一只工作臂可以夹持并旋转瓶盖，符合真实双手分工；
  \item 顶部视野负责看清整张桌面，腕部视野负责看清被机械臂遮挡的瓶口与夹爪附近；
  \item 操作者可以直接示教连续旋转、重新抓取和分类投放，适合形成规模化真实数据。
\end{enumerate}

\figfull[.96\textwidth]{aloha_formal_photo_annotated.png}{项目正式设备照片。照片拍摄时机器人处于静止休眠状态，不是任务执行或成功案例。前侧两只为操作者示教臂，后侧两只为机器人工作臂。}

\section{任务难点不是“抓一次”}

\figfull[.98\textwidth]{task_timeline.pdf}{单个瓶子的七阶段双臂任务链。发现、抓取、找盖、旋拧、分离、投放和回位相互依赖，任一阶段的偏差都可能向后传递；当前阶段完成情况仍由现场观察确认。}

该任务的主要复杂性来自以下耦合：

\begin{itemize}
  \item \textbf{目标不固定：}瓶子可能密集、遮挡、方向变化，也可能没有瓶盖；
  \item \textbf{双臂要同步：}左臂抓得太松，右臂旋转时瓶身会跟转；左臂抓得太紧或位置不当，可能遮挡瓶口；
  \item \textbf{接触具有不确定性：}不同瓶盖的直径、纹理、摩擦和剩余旋紧程度不同；
  \item \textbf{长流程误差累积：}前面一次抓取偏差会影响找盖、旋盖、分离和投放；
  \item \textbf{需要长期循环：}一次成功不代表能够持续工作，回位误差和异常恢复会随时间积累。
\end{itemize}

\section{本报告回答什么}

本报告按照管理层要求，依次回答目的、方法、数据、模型、实验、结果和未来一年计划。报告仅把真实训练记录、保存模型、数据统计、设备照片、训练示教帧、注意力记录和现场观测写入结论；没有录像或计数表的内容一律降级为现场估计。报告不讨论完整三年目标，也不把未来的插管清洗写成已完成能力。
